\clearpage
\begin{landscape}
\chapter*{LAMPIRAN}
%\addcontentsline{toc}{chapter}{LAMPIRAN}	
\section{Pemetaan Topologi dan Infrastruktur}

\subsection{Alokasi Node}
Tabel \ref{tab:topology_mapping} berikut merincikan pemetaan alokasi \textit{node} ke dalam infrastruktur VPS yang digunakan untuk setiap skenario topologi.

{
\scriptsize
\begin{longtable}{|p{2.5cm}|p{3cm}|p{3cm}|p{3cm}|p{3cm}|p{3cm}|p{3cm}|}
\caption{Pemetaan Node ke Infrastruktur VPS} \label{tab:topology_mapping} \\
\hline
\textbf{Topologi} & \textbf{Control Plane} & \textbf{VPS A} & \textbf{VPS B} & \textbf{VPS C} & \textbf{VPS D} & \textbf{VPS E} \\
& (77.237.244.170) & (89.117.50.181) & (185.194.216.132) & (155.133.23.145) & (109.199.99.164) & (109.199.118.220) \\
\hline
\endfirsthead
\hline
\textbf{Topologi} & \textbf{Control Plane} & \textbf{VPS A} & \textbf{VPS B} & \textbf{VPS C} & \textbf{VPS D} & \textbf{VPS E} \\
\hline
\endhead
\hline
\endlastfoot

\textbf{Skenario 1} \newline
PoA: 3 Signer + 1 Non \newline
PoS: 3 Val + 1 Non
& 
Database: PostgreSQL \newline
(Tidak digunakan di akhir)
& 
\textbf{PoA}: Signer-UGM \newline
\textbf{PoS}: Validator-UGM
& 
\textbf{PoA}: Signer-ITB \newline
\textbf{PoS}: Validator-ITB
& 
\textbf{PoA}: Signer-UI \newline
\textbf{PoS}: Validator-UI
& 
\textbf{PoA}: Non-Signer-UNIMED \newline
\textbf{PoS}: Non-Validator-UNIMED
& 
Standby Node \newline
(Sync/Snapshot)
\\ \hline

\textbf{Skenario 2} \newline
PoA: 3 Signer + 3 Non \newline
PoS: 3 Val + 3 Non
& 
\textbf{Monitoring}: \newline
- Ethstats \newline
- Prometheus (SUT/Caliper) \newline
- Grafana (SUT/Caliper)
& 
\textbf{PoA}: Signer-UGM + \newline Non-Signer-UT \newline
\textbf{PoS}: Validator-UGM + \newline Non-Validator-UT
& 
\textbf{PoA}: Signer-ITB + \newline Non-Signer-UNUD \newline
\textbf{PoS}: Validator-ITB + \newline Non-Validator-UNUD
& 
\textbf{PoA}: Signer-UI \newline
\textbf{PoS}: Validator-UI
& 
\textbf{PoA}: Non-Signer-UNIMED \newline
\textbf{PoS}: Non-Validator-UNIMED
& 
Archive/Standby Node
\\ \hline

\textbf{Skenario 3} \newline
PoA: 5 Signer + 5 Non \newline
PoS: 5 Val + 5 Non
& 
\textbf{Orchestration}: \newline
- k3s Kubernetes \newline
- Hyperledger Caliper \newline
- Reports Generator \newline
- Data Analytics
& 
\textbf{PoA}: Signer-UGM + \newline Non-Signer-UT \newline
\textbf{PoS}: Validator-UGM + \newline Non-Validator-UT
& 
\textbf{PoA}: Signer-ITB + \newline Non-Signer-UNUD \newline
\textbf{PoS}: Validator-ITB + \newline Non-Validator-UNUD
& 
\textbf{PoA}: Signer-UI + \newline Non-Signer-Gundar \newline
\textbf{PoS}: Validator-UI + \newline Non-Validator-Gundar
& 
\textbf{PoA}: Signer-UB + \newline Non-Signer-UNIMED \newline
\textbf{PoS}: Validator-UB + \newline Non-Validator-UNIMED
& 
\textbf{PoA}: Signer-ITS + \newline Non-Signer-UNDIP \newline
\textbf{PoS}: Validator-ITS + \newline Non-Validator-UNDIP
\\ \hline

\end{longtable}
}
\end{landscape}

\clearpage
\subsection{Latensi Baseline Jaringan}
Tabel \ref{tab:latency_ping} berikut menunjukkan hasil pengukuran RTT (\textit{Round Trip Time}) menggunakan protokol ICMP dari Control Plane ke masing-masing nodes. Pengukuran dilakukan sebanyak 5 paket per node untuk memastikan konektivitas dasar yang stabil.

\begin{table}[h]
\centering
\caption{Hasil Pengukuran Latensi (Ping) dari Control Plane}
\label{tab:latency_ping}
\begin{tabular}{|l|l|l|c|c|c|c|}
\hline
\textbf{Node} & \textbf{Label Institusi} & \textbf{IP Address} & \textbf{Min (ms)} & \textbf{Avg (ms)} & \textbf{Max (ms)} & \textbf{Mdev (ms)} \\ \hline
VPS A & UGM \& UT & 89.117.50.181 & 0.654 & 1.297 & 3.407 & 1.058 \\ \hline
VPS B & ITB \& UNUD & 185.194.216.132 & 0.536 & 1.243 & 2.071 & 0.664 \\ \hline
VPS C & UI \& Gundar & 155.133.23.145 & 1.042 & 1.804 & 2.695 & 0.638 \\ \hline
VPS D & UB \& UNIMED & 109.199.99.164 & 1.961 & 2.384 & 2.935 & 0.319 \\ \hline
VPS E & ITS \& UNDIP & 109.199.118.220 & 0.780 & 2.345 & 5.576 & 1.738 \\ \hline
\end{tabular}
\end{table}

\subsection{Penggunaan Penyimpanan Blockchain (Disk Usage)}
Berdasarkan pengukuran pada direktori \texttt{chaindata} (data leveldb), berikut adalah penggunaan ruang penyimpanan per \textit{node} di akhir pengujian yang dirincikan pada Tabel \ref{tab:storage_usage}. Terlihat adanya perbedaan ukuran antara \textit{Signer/Validator} (rata-rata 1.6 GB) and \textit{Non-Signer/Non-Validator} (rata-rata 6 GB), yang mengindikasikan beban akumulasi state yang berbeda pada node RPC publik.

\begin{table}[h]
\centering
\caption{Penggunaan Disk Penyimpanan Blockchain (\texttt{chaindata})}
\label{tab:storage_usage}
\begin{tabular}{|l|l|l|c|}
\hline
\textbf{VPS Node} & \textbf{Label Institusi} & \textbf{Peran (Role)} & \textbf{Ukuran Data} \\ \hline
\multirow{2}{*}{VPS A} & UGM & Signer / Validator & 1.7 GB \\ \cline{2-4} 
 & UT & Non-Signer / Sentry & 6.0 GB \\ \hline
\multirow{2}{*}{VPS B} & ITB & Signer / Validator & 1.7 GB \\ \cline{2-4} 
 & UNUD & Non-Signer / Sentry & 6.1 GB \\ \hline
\multirow{2}{*}{VPS C} & UI & Signer / Validator & 1.7 GB \\ \cline{2-4} 
 & Gunadarma & Non-Signer / Sentry & 5.7 GB \\ \hline
\multirow{2}{*}{VPS D} & UB & Signer / Validator & 1.5 GB \\ \cline{2-4} 
 & UNIMED & Non-Signer / Sentry & 5.8 GB \\ \hline
\multirow{2}{*}{VPS E} & ITS & Signer / Validator & 1.5 GB \\ \cline{2-4} 
 & UNDIP & Non-Signer / Sentry & 5.7 GB \\ \hline
\end{tabular}
\end{table}

\clearpage
\subsection{Struktur Direktori Node}
Contoh struktur direktori pada \textbf{VPS A} di akhir pengujian. Struktur ini menunjukkan pemisahan folder antara \textit{Signer} (UGM) dan \textit{Non-Signer} (UT), di mana setiap direktori \texttt{geth} menyimpan data blockchain yang tumbuh seiring beban transaksi.

\begin{small}
\begin{verbatim}
/var/lib/poa
|-- ugm
|   |-- clef
|   |   |-- d6cad8dd3b43fa311d9c
|   |-- geth
|   |   |-- geth
|   |   |   |-- blobpool
|   |   |   |-- chaindata (1.7 GB)
|   |   |   |-- lightchaindata
|   |   |   |-- nodes
|   |   |-- keystore
|-- ut
|   |-- geth
|   |    |-- geth
|   |    |   |-- blobpool
|   |    |   |-- chaindata (6.0 GB)
|   |    |   |-- lightchaindata
|   |    |   |-- nodes
|   |    |-- keystore
\end{verbatim}
\end{small}
